\section{2026 年 11 月核验增补}

本节按论文首次提交时间排列。每项方法附有可核验的外部入选依据。

% METHOD: 2312.15856
\paragraph{SplatMesh：用显式网格组织交互分割和高斯编辑}
2023.12~\cite{2312_15856}。SplatMesh 先从多视图建立带颜色的网格，并按几何和纹理筛减网格，再基于保留的顶点初始化高斯。网格给提示分割和几何操作提供结构，依附网格的高斯继续提高表面渲染质量，用户的编辑提示无需逐条训练模型。

NeRF Synthetic 的视图合成对照与 SuGaR 使用相同细化轮数，另比较了保留三分之一面的网格简化配置。交互编辑示例说明流程用途，不能替代对破碎网格或细结构的普遍准确率评估。原文署名含华为诺亚方舟实验室，库内所列项目链接不作为完整实现依据。

\textbf{入选依据} 论文署名单位包括 Huawei（\href{https://arxiv.org/abs/2312.15856}{论文原文}）。

% METHOD: 2608.26656
\paragraph{CoGeo-GS：联合语义选取、深度补全与外观细化}
2026.08~\cite{2608_26656}。CoGeo-GS 将文本驱动的分割掩码蒸馏到高斯中，为多个目标分配语义身份。移除对象后，以对齐尺度的单目深度约束扩散补全，再用几何正则化外观细化缓和跨视图与边界不一致。

作者在 Mip-NeRF 360 和 SPIn-NeRF 上分别验证单对象和多对象移除，使用给定相机内参与 COLMAP 点云初始化，并与 InFusion、SPIn-NeRF、GaussianEditor 和 Gaussian Grouping 比较。被遮挡区域只能由条件补全，论文的视觉与重建指标针对该数据和相机设置。

\textbf{入选依据} 论文署名单位包括 ByteDance（\href{https://arxiv.org/abs/2608.26656}{论文原文}）。

% METHOD: 2609.03341
\paragraph{PointGT：让点表示在形变后仍能读取编辑纹理}
2026.09~\cite{2609_03341}。PointGT 在神经点渲染器 PAPR 上学习三维点和二维纹理图之间的映射，以两种几何正则约束表面。由于形变会改变射线交点位置，方法用位移融合把变形后的交点映射回规范空间，让编辑过的纹理继续贴合表面。

论文比较纹理和几何编辑及多视角渲染，复杂拓扑处仍难得到低畸变纹理图。Pratul P. Srinivasan 的 Google DeepMind 署名可核实。PointGT 使用无空间尺度的神经点和注意力渲染，适合作为高斯编辑的对照路线。

\textbf{入选依据} 论文署名单位包括 Google DeepMind（\href{https://arxiv.org/html/2609.03341v1}{论文原文}）。作为点渲染编辑的路线对照，核心表示采用 PAPR。

% METHOD: 2609.03349
\paragraph{P-CORE：以自监督形变改善神经点表面连续性}
2026.09~\cite{2609_03349}。P-CORE 在 PAPR 的点注意力中提取预测表面点，并给原始点施加已知解析形变。用同一形变作用于规范表面点，再与对变形场景重新查询的表面作比较，配合纹理一致性约束细化模型。训练无需实际采集变形后的场景。

作者在 Neural Editor 与 Objaverse 以相同形变比较渲染和几何伪影。过度细调会平滑外观。作者仓库确有训练、微调和测试入口，但算法采用无空间尺度的 PAPR 注意力点，不能当成 3DGS 的一项形变基元方法。

\textbf{入选依据} \href{https://github.com/zvict/p-core}{公开实现}已找到与论文对应的代码文件。作为点渲染编辑的路线对照，核心表示采用 PAPR。
